\section{Exact linearization, transfer functions, and channels}

We consider the validated strong-Allee predator--prey model with parameters
\(r=3/2\), \(K=1\), \(a=1\), \(h=1/2\), \(e=4/5\), \(m=2/5\)~\cite[Sec.~4]{source_pack/04_EXACT_STRONG_ALLEE_BASELINE.md}.
The dynamics are
\begin{align}
\dot x &= r x\Bigl(1-\frac{x}{K}\Bigr)\Bigl(\frac{x}{A}-1\Bigr)-\frac{a x y}{1+h x},\\
\dot y &= e\,\frac{a x y}{1+h x}-m y,
\end{align}
where \(A\in(0,2/3)\) is the Allee threshold.

\subsection{Equilibrium and Jacobian}
The coexistence equilibrium \((x^*,y^*)\) is given by
\begin{equation}
\label{eq:coex}
x^*=\frac{m}{a(e-mh)}=\frac23,\qquad
y^*=\frac{1+hx^*}{a}\,r\Bigl(1-\frac{x^*}{K}\Bigr)\Bigl(\frac{x^*}{A}-1\Bigr)
      =\frac{2(2-3A)}{9A}.
\end{equation}
The Jacobian at coexistence is
\begin{equation}
\label{eq:jac}
J(A)=\begin{pmatrix}
\dfrac{7A-2}{8A} & -\dfrac12 \\[1ex]
\dfrac{2-3A}{10A} & 0
\end{pmatrix},
\end{equation}
with trace, determinant, and discriminant
\begin{align}
T(A)&=\operatorname{tr}J(A)=\frac{7A-2}{8A},\label{eq:trace}\\
D(A)&=\det J(A)=\frac{2-3A}{20A},\label{eq:det}\\
T(A)^2-4D(A)&=\frac{(19A-10)(23A-2)}{320A^2}.\label{eq:disc}
\end{align}
The Matignon stability bound for a commensurate Caputo system of order
\(0<\alpha<1\) is
\begin{equation}
\label{eq:alpha}
\alpha^*(A)=\frac{2}{\pi}\arctan\!\Bigl(
\frac{\sqrt{4D(A)-T(A)^2}}{T(A)}\Bigr),
\end{equation}
and the coexistence equilibrium is locally asymptotically stable iff
\(0<\alpha<\alpha^*(A)\).  For the stress-test baseline \(A=3/10\),
\(\alpha^*(3/10)\approx0.9690122761517084\).

\subsection{Controllability and observability}
The pair \((J,B)\) is controllable for either prey actuation
\(B=e_1\) or predator actuation \(B=e_2\), because
\(\det\mathcal C_x = \det[B,JB] = \frac{2-3A}{10A}>0\) and
\(\det\mathcal C_y = \det[B,JB] = \frac12>0\).  
Similarly, the pair \((C,J)\) is observable for either prey-only
measurement \(C=e_1^\top\) or predator-only measurement
\(C=e_2^\top\), since \(\det\mathcal O_x = -\frac12\neq0\) and
\(\det\mathcal O_y = -\frac{2-3A}{10A}\neq0\).  
Thus a single‑species actuation and sensing suffice to excite and
observe both linearized modes.

\subsection{Transfer functions}
Introduce the fractional temporal symbol
\(q_\alpha(s)=\tau_0^{-1}(s\tau_0)^\alpha\) with \(\tau_0=1\) (without loss of
generality) and the characteristic polynomial
\(\Delta_\alpha(s)=q_\alpha^2 - T q_\alpha + D\).
The resolvent is
\[
(q_\alpha I-J)^{-1}
=\frac{1}{\Delta_\alpha(s)}
\begin{pmatrix}
q_\alpha & \frac12 \\[0.5ex]
\frac{2-3A}{10A} & q_\alpha-T
\end{pmatrix}.
\]
Hence the four scalar transfer channels are
\begin{align}
G_{x\leftarrow x}(s)&=\frac{q_\alpha}{\Delta_\alpha(s)},\label{eq:Gxx}\\
G_{y\leftarrow x}(s)&=\frac{2-3A}{10A}\,\frac{1}{\Delta_\alpha(s)},\label{eq:Gyx}\\
G_{x\leftarrow y}(s)&=-\frac12\,\frac{1}{\Delta_\alpha(s)},\label{eq:Gxy}\\
G_{y\leftarrow y}(s)&=\frac{q_\alpha-T}{\Delta_\alpha(s)}.\label{eq:Gyy}
\end{align}
For the collocated channels (\(CB=1\)) the high‑frequency expansion gives
\(G(s)\sim q_\alpha^{-1}\sim s^{-\alpha}\), whereas the cross‑species
channels (\(CB=0\)) decay as \(s^{-2\alpha}\) or faster.  
Thus a collocated perturb‑and‑measure protocol yields the cleanest
high‑frequency signature of the fractional order \(\alpha\).

\end{document}